\documentclass[10pt,a4paper]{amsart}
\usepackage[margin=19mm]{geometry}
\usepackage{amsmath,amssymb,mathtools}
\usepackage[scr=rsfs,cal=euler]{mathalfa}
\usepackage{xfrac}
\usepackage[unicode,hidelinks,bookmarks=false]{hyperref}
\newcommand{\ie}{i.e.\ }
\newcommand{\cf}{cf.\ }
\newcommand{\resp}{respectively\ }
\newcommand{\Subsection}{\subsection}
\providecommand{\nobreakdash}{\nobreak\hbox{-}}
\setlength{\emergencystretch}{2em}
\multlinegap0pt
\usepackage{amssymb}
\usepackage{enumerate}
\usepackage{xcolor}
\newtheorem{lem}{Lemma}
\newtheorem{prop}{Proposition}
\usepackage{cleveref}
\crefname{lem}{Lemma}{Lemmas}
\crefname{prop}{Proposition}{Propositions}
\renewcommand{\bmod}{\,\,\mathsf{mod}\,\,}
\newcommand{\car}{\mathsf{char}}
\renewcommand{\geq}{\geqslant}
\renewcommand{\leq}{\leqslant}
\newcommand{\mc}[1]{\mathcal{#1}}
\newcommand{\mfm}{\mathfrak{m}}
\newcommand{\mg}[1]{#1^{\times}}
\renewcommand{\min}{\mathsf{min}}
\newcommand{\nat}{\mathbb{N}}
\DeclareMathOperator{\newsum}{\mathsf{\Sigma}}
\renewcommand{\setminus}{\smallsetminus}
\newcommand{\sos}[2]{\sum_{#1}\!{#2}^2}
\renewcommand{\sum}{\newsum}
\newcommand{\sums}[1]{\sum\!{#1}^2}
\title{Sums of squares in function fields over a henselian valued field}
\author{Karim Johannes Becher, Nicolas Daans, David Grimm, Gonzalo Manzano Flores, Marco Zaninelli}
\date{Source: arXiv:2302.11425v4 (CC BY 4.0)}
\begin{document}
\maketitle

\noindent\emph{Source and licence.} Sums of squares in function fields over a henselian valued field. Version arXiv:2302.11425v4 (CC BY 4.0), distributed by arXiv under the Creative Commons Attribution 4.0 International licence, http://creativecommons.org/licenses/by/4.0/, as recorded in the arXiv licence metadata for this exact version and shown on the abstract page https://arxiv.org/abs/2302.11425v4. Full licence notice: \texttt{licenses/arxiv-2302.11425-CC-BY-4.0.txt}.\par\smallskip

\noindent\textit{Editorial scope.} Counted unit: the article's prop p' (source0.tex lines 346--351). The article's complete original proof of it is delivered with it (source0.tex lines 353--373, one \texttt{proof} environment of 21 lines). No mathematics is written, repaired or re-derived: the statement and the proof are the article's own lines, in the article's own notation, and no retained line is substituted or re-flowed.  Retained uncounted as the source-local context the proof needs: lines 331--336.  Nothing was omitted from any retained span, so the package carries no omission record.  No retained line contains a citation, so the wrapper carries no bibliography and no bibliography entry.  No retained line contains a figure environment, an image-inclusion command or drawing code, so no image is delivered and the package declares no asset dependency.  Editorial formatting only, in addition: an AMS \texttt{amsart} preamble that restates the article's own declarations and macros used by the retained lines, namely the environments \texttt{lem}, \texttt{prop} on separate counters, together with the standard packages those lines need.  The retained environments print in this document as Lemma 1, Proposition 1 in order of appearance, whereas the article prints the same items with the numbering of its own full document.  The complete native source of the article as distributed in the arXiv e-print for this exact version is bundled unchanged as \texttt{sources/137729/source0.tex}.
\begin{lem}\label{L:valuation-residue-level}
Let $v$ be a valuation on $K$ and 
 $n\in\nat$ such that $1\leq n\leq s(Kv)$.
Then $v(x_1^2+\dots+x_n^2)=2\min\{v(x_i)\mid 1\leq i\leq n\}$ for any $x_1,\dots,x_n\in K$.
In particular $\mc{O}_v\cap\sos{n}{K}=\sos{n}{\mc{O}}_v$.
\end{lem}

\begin{prop}\label{inequality p'}
Let $v$ be a valuation on $K$.
We have that  $s(K)\geq s(Kv)$ and $p(K)\geq p(Kv)$.  
Furthermore, if $\car(Kv)\neq 2$ and $v$ is henselian, then 
\[s(K)=s(Kv)\quad \mbox{ and }\quad p'(K)= p'(Kv).\]
\end{prop}

\begin{proof}
It follows from \Cref{L:valuation-residue-level} that, for any natural number $n\leq s(Kv)$ and any $x\in\mg{\mc{O}}_v\cap\sos{n}{K}$, one has $x+\mfm_v\in\sos{n}{(Kv)}$.
Using this and the fact that $p(Kv)\leq s(Kv)+1$, one readily sees that $s(K)\geq s(Kv)$ and $p(K)\geq p(Kv)$. 

Assume now that $v$ is henselian and $\car(Kv)\neq 2$.

We consider first the case where $s(Kv)<\infty$.
We set $s=s(Kv)$ and choose $x_1,\dots,x_s\in {\mc O}_v$ such that $x_1^2+\dots+x_s^2 \equiv -1\bmod \mfm_v$. 
Since $v$ is henselian and $\car(Kv)\neq 2$, we can  find $y\in \mc{O}_v$ such that $1+x_1^2+\dots+x_{s-1}^2+y^2 =0$. Hence $s(K)=s=s(Kv)$ and consequently $p'(K)=s+1=p'(Kv)$.

We now consider the case where $s(Kv)=\infty$. Then $s(K)=\infty=s(Kv)$ and in particular $p'(K)=p(K)$ and $p'(Kv)=p(Kv)$. Since $p(K)\geq p(Kv)$, to show equality, we may assume that $p(Kv)<\infty$. Let $r=p(Kv)$.
Consider $x\in \sums{K}\setminus\{0\}$.
Then $v(x)\in 2 vK$, by \Cref{L:valuation-residue-level}, so there exists $t\in \mg K$ such that $t^2x\in\mg{\mc{O}}_v$.
In view of \Cref{L:valuation-residue-level}, it follows that there exist $x_1,\dots,x_{r-1}\in \mc{O}_v$ and $x_r\in\mg{\mc{O}}_v$ such that 
$t^2x\equiv x_1^2+\dots+x_r^2\bmod \mfm_v$.
Since $v$ is henselian and $\car(Kv)\neq 2$, we obtain that 
$t^2x=x_1^2+\dots+x_{r-1}^2+y^2$ for some $y\in \mc{O}_v$. Hence $t^2x\in \sos{r}{K}$ and therefore $x\in\sos{r}{K}$.
This argument shows that $p(K)=p(Kv)$. 

Therefore $p'(K)=p'(Kv)$ holds in every case.
\end{proof}

\end{document}
